
\begin{proposition}
\label{V.6.4}
Sea~$X$ un objeto punteado conexo de~$\cal C$, asociado a un subgrupo
abierto~$U$ de~$\pi_F$. Para que
\ifthenelse{\boolean{orig}}{$u$}{$U$}
contenga a~$u(\pi_{F'})$ (\resp al subgrupo normal cerrado engendrado
por~$u(\pi_{F'})$), es necesario y suficiente que~$H(X)$ admita una
sección punteada (\resp esté completamente descompuesto).
\end{proposition}

Se llamará \emph{sección} ---sobrentendido: sobre el objeto final--- de un
objeto~$X$ de una categoría galoisiana~$\cal C$ a un morfismo del objeto
final~$e_{\cal C}$ en~$X$, lo que equivale a dar un elemento~$a$ de~$F(X)$
invariante por~$\pi_F$; si~$X$ es punteado, se dice que se tiene una
\emph{sección punteada} si es compatible con las estructuras punteadas de
$X$ y~$e_{\cal C}$, \ie si~$a$ es precisamente el
\ifthenelse{\boolean{orig}}
{objeto marcado}
{elemento marcado}
de~$F(X)$. Tal sección es, pues, única, y existe si y solo si el
\ifthenelse{\boolean{orig}}
{objeto marcado}
{elemento marcado}
de~$F(X)$ es invariante por~$\pi_F$. Por último, se dice que un objeto de
una categoría galoisiana está \emph{completamente descompuesto}
\index{completamente descompuesto (objeto)|hyperpage}%
\index{descompuesto (objeto completamente)|hyperpage}%
si es isomorfo a una suma de objetos finales, \ie si~$\pi_F$ actúa
trivialmente sobre~$F(X)$, condición evidentemente más fuerte que la
existencia de una sección punteada cuando~$X$ es punteado. La
Proposición~\Ref{V.6.4} resulta inmediatamente de las definiciones y
observaciones anteriores.

\begin{corollaire}
\label{V.6.5}
Para
\marginpar{138}
que~$u$ sea trivial, es necesario y suficiente que, para todo objeto~$X$
de~$\cal C$, $H(X)$ esté completamente descompuesto.
\end{corollaire}

\begin{proposition}
\label{V.6.6}
Sea~$X'$ un objeto punteado conexo de~$\cal C'$, asociado a un subgrupo
abierto~$U'$ de~$\pi_{F'}$. Para que~$U'$ contenga a~$\Ker u$, es necesario
y suficiente que exista un objeto punteado conexo~$X$ de~$\cal C$ y un
homomorfismo punteado de la componente conexa punteada~$X'_0$ de~$H(X)$
en~$X'$ (por tanto, que
\ifthenelse{\boolean{orig}}{$X$}{$X'$}
sea isomorfo, como objeto punteado, a un cociente de la componente conexa
neutra de la imagen inversa de un objeto punteado de~$\cal C$). Si~$u$ es
sobreyectivo, la condición anterior equivale también a la siguiente:
\ifthenelse{\boolean{orig}}{$X$}{$X'$}
es isomorfo a un objeto de la forma~$H(X)$, donde~$X$ es un objeto punteado
de~$\cal C$.
\end{proposition}

(Se llama \emph{componente conexa neutra}
\index{neutra (componente conexa de un objeto punteado)|hyperpage}%
de un objeto punteado~$X$ de una categoría galoisiana~$\cal C$ al único
subobjeto conexo punteado de~$X$; en virtud de~\Ref{V.5.3}, corresponde a
la órbita bajo~$\pi_F$ del punto marcado de~$F(X)$). Como el hecho de que
$U'$ contenga a~$\Ker u$ no depende del punto elegido en~$X'$ (pues elegir
otro punto equivale a reemplazar
\ifthenelse{\boolean{orig}}{$U$}{$U'$}
por un subgrupo conjugado), se obtiene:

\begin{corollaire}
\label{V.6.7}
Para que~$U'$ contenga a~$\Ker u$, es necesario y suficiente que exista un
objeto~$X$ de~$\cal C$ (que puede suponerse conexo) y un morfismo de una
componente conexa de~$H(X)$ en~$X'$. Si~$u$ es sobreyectivo, esto significa
también que~$X'$ es isomorfo a un objeto de la forma~$H(X)$.
\end{corollaire}

\begin{corollaire}
\label{V.6.8}
Para que~$u$ sea inyectivo, es necesario y suficiente que, para todo
objeto~$X'$ de~$\cal C'$, exista un objeto~$X$ de~$\cal C$ y un
homomorfismo de una componente conexa de~$H(X)$ en~$X'$.
\end{corollaire}

\begin{proposition}
\label{V.6.9}
Las condiciones siguientes son equivalentes:
\begin{enumerate}
\item[(i)] El homomorfismo $u\colon\pi_{F'}\to\pi_F$ es sobreyectivo.
\item[(ii)] Para todo objeto conexo~$X$ de~$\cal C$, $H(X)$ es conexo.
\item[(iii)] El funtor~$H$ es plenamente fiel.
\end{enumerate}
\end{proposition}

Esta última condición significa que, para dos objetos~$X$, $Y$ de~$\cal C$,
la aplicación natural
$$
\Hom(X,Y)\to\Hom(H(X),H(Y))
$$
es biyectiva.

\begin{corollaire}
\label{V.6.10}
Para
\marginpar{139}
que~$u$ sea un isomorfismo, es necesario y suficiente que~$H$ sea una
equivalencia de categorías, o, equivalentemente, que se verifiquen las dos
condiciones siguientes:
\begin{enumerate}
\item[a)] para todo objeto conexo~$X$ de~$\cal C$, $H(X)$ es conexo;
\item[b)] todo objeto de~$\cal C'$ es isomorfo a un objeto de la
forma~$H(X)$.
\end{enumerate}
\end{corollaire}

\begin{proposition}
\label{V.6.11}
Sean $H\colon\cal C\to\cal C'$ y
\ifthenelse{\boolean{orig}}
{$H'\colon\cal C\to\cal C''$}
{$H'\colon\cal C'\to\cal C''$}
funtores exactos entre categorías galoisianas; sea~$F''$ un funtor
fundamental sobre~$\cal C''$, pongamos $F'=F''H'$ y $F=F'H$, y consideremos
los homomorfismos asociados
$$
u'\colon\pi_{F''}\to\pi_{F'}\qquad
u\colon\pi_{F'}\to\pi_F.
$$
Para que
\ifthenelse{\boolean{orig}}
{$\Ker u\subset\Im u'$}
{$\Ker u\supset\Im u'$},
\ie para que~$uu'$ sea el homomorfismo trivial, es necesario y suficiente
que, para todo objeto~$X$ de~$\cal C$, $H'(H(X))$ esté completamente
descompuesto. Para que
\ifthenelse{\boolean{orig}}
{$\Ker u\supset\Im u'$}
{$\Ker u\subset\Im u'$},
es necesario y suficiente que, para todo objeto punteado conexo~$X'$ de
$\cal C'$ tal que~$H'(X')$ admita una sección punteada, exista un objeto~$X$
de~$\cal C$ y un homomorfismo de una componente conexa de~$H(X)$ en~$X'$.
\end{proposition}

La primera afirmación resulta de la última afirmación de~\Ref{V.6.4}. La
segunda resulta de la conjunción de~\Ref{V.6.4} y~\Ref{V.6.6}.

\begin{remarque}
\label{V.6.12}
En general, bajo las condiciones de~\Ref{V.6.8}, no es cierto que~$X'$ sea
isomorfo a un objeto de la forma~$H(X)$. Puede demostrarse que, para que
todo objeto conexo (y, por tanto, todo objeto) de~$\cal C'$ sea isomorfo a
un objeto de la forma~$H(X)$, es necesario y suficiente que~$u$ sea un
isomorfismo de~$\pi_{F'}$ sobre un subgrupo \emph{factor directo}
de~$\pi_F$. Sin embargo, en la práctica se construye directamente un
homomorfismo $\pi_F\to\pi_{F'}$
\ifthenelse{\boolean{orig}}
{inverso por la derecha}
{inverso por la izquierda}
de~$u$ mediante un funtor exacto conveniente de~$\cal C'$ en~$\cal C$.
\end{remarque}

\begin{proposition}
\label{V.6.13}
Sean~$\cal C$ una categoría galoisiana provista de un funtor fundamental~$F$,
$S$ un objeto conexo de~$\cal C$ y~$\cal C'$ la categoría de los objetos
de~$\cal C$ sobre~$S$. Entonces~$\cal C'$ es una categoría galoisiana, y el
funtor $X\mto H(X)=X\times S$ de~$\cal C$ en~$\cal C'$ es exacto. Sea
$a\in F(S)$, y sea~$F'$ el funtor de~$\cal C'$ en la categoría de los
conjuntos finitos definido por
$$
F'(X')=\text{preimagen de~$a$ por }F(X')\to F(S).
$$
Entonces se tiene un isomorfismo $F\cong F'\circ H$, y el homomorfismo
correspondiente
$$
u\colon\pi_{F'}\to\pi_F
$$
\marginpar{140}%
es un isomorfismo de~$\pi_{F'}$ sobre el subgrupo abierto~$U$ de~$\pi_F$
estabilizador del elemento marcado~$a$ de
\ifthenelse{\boolean{orig}}{$F(X)$}{$F(S)$}.
\end{proposition}

La demostración se deja al lector.
